\cotefile{15}
\scanpage{34}
% archive foliation: 33

\uline{Ex. 3} Soit fixé, en plus de la catégorie de Deuring $C$ sur $A$ comme dans Ex. 2, un élément $\tau$ de $V_A(C)$. Si $E$ est \inserted{\unsure{centrale}} de rang $n$, il faut considérer que les $C$ qui \uline{\unsure{réalisent}} les
\deleted{(\unsure{de manière à convenir à la rigueur})}
\unsure{gerbes} de lien $G_m$ \uuline{sur $A$} trivialisées par $n$,
\deleted{et \unsure{sur} $A$ (car si $E$ est \unsure{définie} sur $A$ \unsure{complète de} \unsure{$G_m$})}
\inserted{$C$ qui correspond au \unsure{corps} \deleted{\unsure{grp}}, \unsure{dans $K$}, cela revient à une algèbre d'Azumaya sur $R$, i.e. la classe de $E$ dans $\mathrm{Br}(K)$ est dans l'image de $\mathrm{Br}(R)$.}
et il faut voir de \unsure{même} que les couples $(C,\tau)$ \unsure{équivalent} à la donnée d'une gerbe de groupe $\mu_n$.

\margin{\unsure{et trivialisation}\\
\unsure{gerbes} sur $A$ de lien $G_m$ $\longrightarrow$ catégorie $C$ des \unsure{représentations} \unsure{faisant} \unsure{des} \unsure{objets} \unsure{fidèles} \unsure{vis-à-vis} du groupe, poids $1$.}
\margin{$C\longmapsto$ \unsure{gerbes} des foncteurs admissibles de $C$ dans les \unsure{modules} \unsure{``caractéristiques''} de $C$\\
cat. de Deuring sur \unsure{$E$}\\
\unsure{$\longleftarrow$} \unsure{$\xi\in\widehat{\mathcal G}$}\\
\unsure{$\widehat{\mathcal G}=[\mathrm{Br}\,A\longrightarrow\mathrm{Br}\,K]$}}

Pour cela, \deleted{\unsure{on}} introduire \deleted{\unsure{la catégorie des gerbes}} pour une gerbe de lien $\mu_n$ sur $A$ \unsure{définir} bien une $(C,\tau)$:

$C$: les \deleted{\unsure{trivialisations}} \deleted{de} \unsure{représentations} de degré $n$, de rang $1$ dans la gerbe, $\tau=$ \deleted{\unsure{la représentation déterminant}}.

Nous allons déterminer d'abord \unsure{quels sont} les autoéquivalences de $(E,\tau)$; en termes d'un objet $X$ de $C$, \inserted{d'où $O=\operatorname{End}(X)$,} c'est la catégorie des \deleted{\unsure{algèbres}} bimodules $P$ \unsure{sous} $O$, inv. loc. sur $A$ des deux côtés, munis d'un isom
\[
N_{O/A}(P_{\unsure{O|\sigma}})\simeq\deleted{\unsure{O}}A.
\]
\inserted{dans $P(A)$}
\margin{(\deleted{$\unsure{\nu}(X)=L\otimes$} $\tau=L\otimes_A\nu(X)$)}
\margin{\[
\begin{aligned}
\varphi(X)&=P\otimes_O X\\
\nu(\varphi(X))&=N(P)\otimes\nu(X)\\
&\unsure{\simeq}\\[-6pt]
&\varphi(\nu(X))
\end{aligned}
\]}
\deleted{\unsure{compatible avec}}

\deleted{Supposons que tous les $P_{\unsure{O|\sigma}}$ \inserted{\unsure{de} $\unsure{P(O/A)}$} soient isomorphes, ce qui est le cas si $A$ est local, \unsure{ou} semi local.}

\unsure{Alors} c'est une Gr.-catégorie, \inserted{et \unsure{même} une catégorie de Picard \unsure{str. comm.}} \unsure{où} les endom. de l'objet $1$ sont les éléments $\lambda\in A^*$ tels que $\lambda^n=1$; \unsure{cette} catégorie de Picard contient comme sous-catégorie pleine $\uline{\operatorname{Tors}}(S,\mu_n)$ des
